\Needspace{8\baselineskip}\originalchapter{2006 作业六: Alfv\'en 波与 sausage}
题源: 2 Nov 2006, due 9 Nov, 原件包括提示页.
\begin{exercise}\label{ex:06-6-1}
原题1. 弦波$\omega=k\sqrt{\mathcal T/\lambda_m}$, $\mathcal T$为张力、$\lambda_m$为线质量. (a)以圆磁通管类比给Alfv\'en色散; (b)朴素结果差\ensuremath{\sqrt2}, 因忽略侧磁压. 管位移 $y=y_0\sin(kz-\omega t)$, a\ensuremath{\ll}2\ensuremath{\pi}/k、ky0\ensuremath{\ll}1; 给侧压回复每长度 $B^2\pi a^2/(2\mu_0R)$, 求恢复力并补正确速度.
\end{exercise}
\begin{solution}
(a)仅端面Maxwell张力 $T_B=B^2A/(2\mu_0)$, 线质量\ensuremath{\rho}A, 朴素\ensuremath{\omega}\ensuremath{{}^{2}}=k\ensuremath{{}^{2}}B\ensuremath{{}^{2}}/(2\ensuremath{\mu}0\ensuremath{\rho}). (b)小曲率 $1/R\simeq-y''$取回复方向, 侧力同为 $(B^2A/(2\mu_0))y''$. 总力 $B^2Ay''/\mu_0$, $\rho A\ddot y=B^2Ay''/\mu_0$, 得\ensuremath{\omega}=kVA. 截面面积抵消, 合并端面与周边应力才等于流体张力恒等式.
\end{solution}
\begin{exercise}\label{ex:06-6-2}
原题2. 均匀无流无电流背景, 理想MHD只将Ohm改为E+v\ensuremath{\times}B=\ensuremath{\eta}J. (a)完整线性方程; (b) $B_1=-\mathbf k\times(\mathbf v_1\times\mathbf B_0)/\omega^*$求\ensuremath{\omega}*; (c)理想速度方程含cA与cs, 求替代cA*、cs*; (d)平行剪切Alfv\'en色散; (e)驱动\ensuremath{\omega}实, k=kr+iki求二者.
\end{exercise}
\begin{solution}
(a) $\partial_t\rho_1+\rho_0\nabla\cdot v_1=0$, $\rho_0\partial_tv_1=J_1\times B_0-\nabla p_1$, $p_1=c_s^2\rho_1$, $\mu_0J_1=\nabla\times B_1$, divB1=0, $E_1+v_1\times B_0=\eta J_1$, $\partial_tB_1=-\nabla\times E_1$. (b) $\omega^*=\omega+\ii\eta k^2/\mu_0$. (c) $c_A^{*2}=c_A^2\omega/\omega^*$, $c_s^*=c_s$. 磁项被扩散改变, 气压闭合不变. (d) $\omega(\omega+\ii D_mk^2)=k^2V_A^2$. (e) $k^2=\omega^2/(V_A^2-\ii\omega D_m)$, 定义\ensuremath{\theta}=atan(\ensuremath{\omega}Dm/VA\ensuremath{{}^{2}}),
\begin{equation}
 k_r=\frac{\omega\cos(\theta/2)}{(V_A^4+\omega^2D_m^2)^{1/4}},\quad
 k_i=\frac{\omega\sin(\theta/2)}{(V_A^4+\omega^2D_m^2)^{1/4}}.
\end{equation}
\ensuremath{\omega}>0选ki>0, 空间因子e$^{ikz}$向+z衰减. 弱电阻 $k_i\simeq\omega^2D_m/(2V_A^3)$, 与时间阻尼不是同一单位量.
\end{solution}
\begin{exercise}\label{ex:06-6-3}
原题3. 直z pinch无流, $(p+B_\theta^2/(2\mu_0))'+B_\theta^2/(\mu_0r)=0$, 不可压扰动 $\mathbf V(r)\e^{\ii kz+\ii\omega t}$, Vr(a)=0. (a)由 $B_1=\nabla\times(V\times B)/(\ii\omega)$求B1\ensuremath{\theta}; (b)求r/z动量; (c)合不可压条件化单二阶Vr方程; (d)乘rVr积分证明B\ensuremath{\theta}\ensuremath{{}^{\prime}}\ensuremath{-}B\ensuremath{\theta}/r>0全区时失稳; (e)等价于 $2B_\theta^2/r+\mu_0p'<0$, 解释仅粒子漂移的``好坏曲率''判断为何不足.
\end{exercise}
\begin{solution}
(a) $B_{1\theta}=(B/r-B')V_r/(\ii\omega)$. (b)令 $H=B B_{1\theta}+\mu_0p_1$, 则 $\ii\omega\rho\mu_0V_r=-2BB_{1\theta}/r-H'$, $\ii\omega\rho\mu_0V_z=-\ii kH$. (c)不可压 $(rV_r)'/r+\ii kV_z=0$, 消H得到
\begin{equation}
 \left[\frac\rho r(rV_r)'\right]'
 -k^2\left[\rho+\frac{2B(B'-B/r)}{\mu_0r\omega^2}\right]V_r=0.
\end{equation}
(d)轴正规、壁Vr0使边界积分零, 对实径向试函数
\begin{equation}
 \omega^2=-\frac{(2k^2/\mu_0)\int_0^a B(B'-B/r)V_r^2\dd r}
 {\int_0^a(\rho/r)[(rV_r)']^2\dd r+k^2\int_0^a\rho rV_r^2\dd r}.
\end{equation}
若B>0且B\ensuremath{{}^{\prime}}\ensuremath{-}B/r>0全区, 分母正分子负, \ensuremath{\omega}\ensuremath{{}^{2}}<0有增长支. (e)由平衡 $2B^2/r+\mu_0p'=B^2/r-BB'=-B(B'-B/r)$, 等价. 曲率既给压力驱动也给磁恢复, 只看粒子中心漂移方向遗漏磁能; 此不等式为足够的全区条件, 局部单点变号不直接推出同一全局结论.
\end{solution}
